\section{The undivided energy and its differentiation}
\label{sec:energy}

The energy density of a field point, $\tfrac12(|\bB|^2+|\bE|^2)$, is a single quantity: it treats the magnetic and the electric triple
alike and it is unchanged by duality. A field point also has a rest energy, the energy density seen in a frame in which its Poynting vector
vanishes. This section asks how the second is contained in the first: which quadratic energy the arena singles out and what its symmetry is,
how that energy differentiates into a Lorentz-invariant rest part and a part carried by motion without any interaction or second field, where
the conversion constant $c$ enters, and which of the arena's own structures can carry the name of a mass. The relation $E=mc^2$
\cite{Einstein1905} is inherited; what is at stake is only how the arena realises its ingredients. Most of the answers are inherited
electrodynamics written in the arena's variables, and are said to be so; the arena adds which of its forms plays which role. Throughout, $c$
enters only as the conversion between the two triples of \cref{thm:onec}, never as a symmetry.

\paragraph{Notation.} We use the field-strength reading of \cref{thm:emfield}: $X=\sum_i(B_ie_i+E_ie_{i+3})$, $z=\bB+i\bE\in\C^3$, the
binding root $G$ acts as multiplication by $i$, and $q(z)=z\cdot z$ has $\Re q=\eta(X,X)=|\bB|^2-|\bE|^2$ and $\Im q=2\,\bE\cdot\bB=-\eta(X,GX)$.
With the Hermitian product $h(z,w)=\sum_j\bar z_jw_j$ we put
\begin{equation}
\label{en:defs}
\begin{gathered}
E_6:=\tfrac12|z|^2=\tfrac12\bigl(|\bB|^2+|\bE|^2\bigr),\qquad E_0:=\tfrac12|q(z)|,\\
\mathbf S:=\bE\times\bB=\tfrac i2\,\bar z\times z=-\mathbf S_{\rm t},\qquad \mathbf S_{\rm t}:=\bE\times\Btext .
\end{gathered}
\end{equation}
$E_6$ is the energy density $u$ of \cref{sec:orth}, $\mathbf S$ the Poynting vector of the frame (\cref{ex:circular}) and $\mathbf S_{\rm t}$ the
textbook one. Applied to a frame field $X$ these are the frame field's energies; for the momentum field $\pi=\eta^{-1}p$ of the interior dynamics studied below, \eqref{eq:splitL}
and \cref{sec:closure} we write $E_6(\pi)$, $E_0(\pi)$ and $\mathbf S(\pi)$. Units are Heaviside--Lorentz with $c=1$ until
\cref{sec:energy-c}. The generators $L_i,K_i$ of $\lG$ (\cref{thm:emfield}) act by $\delta\bB=e_i\times\bB$, $\delta\bE=e_i\times\bE$ and by
$\delta\bE=e_i\times\bB$, $\delta\bB=-e_i\times\bE$ respectively, so that $\boldsymbol\omega\cdot L+\bv\cdot K$ acts on $z$ as
$z\mapsto(\boldsymbol\omega+i\bv)\times z$. We write $\mathfrak p_G$ for the part of $\so{3,3}$ anticommuting with $G$, so that
$\so{3,3}=\lG\oplus\mathfrak p_G$, and the Lorentz-class inertia of \cref{thm:inertiareq} as $N=(\nu_1\eta^{-1}+\nu_2G\eta^{-1})/m$, its
coefficients renamed so that the letter $b$ is free for a gauge generator in \cref{sec:energy-transport}.

\subsection{\texorpdfstring{One undivided energy, with pointwise symmetry $U(3)$}{One undivided energy, with pointwise symmetry U(3)}}
\label{sec:energy-unique}

\begin{theorem}[Uniqueness of the undivided energy]
\label{en:unique}
\begin{enumerate}[label=(\roman*),leftmargin=2.2em]
\item The symmetric forms on $V$ invariant under the diagonal rotations are spanned by $|\bB|^2$, $|\bE|^2$ and $\bB\cdot\bE$, and they split as
$\R\,\id\oplus\operatorname{span}\{\eta,\eta G\}$, the first summand commuting and the second anticommuting with $G$.
\item A quadratic energy invariant under the diagonal rotations and under the duality rotations $e^{\vartheta G}$ is a multiple of $E_6$, which is
positive definite.
\item The forms invariant under all of $\lG$ are $\operatorname{span}\{\eta,\eta G\}$, that is $|\bB|^2-|\bE|^2$ and $\bE\cdot\bB$; none of them is
definite.
\end{enumerate}
\end{theorem}

\begin{proof}
The diagonal rotations act on $V=V^+\oplus V^-$ as $\R^2\otimes\R^3$ with $SO(3)$ on the second factor, which is absolutely irreducible, so by
Schur's lemma \cite{GoodmanWallach} the invariant forms are $N=\bigl(\begin{smallmatrix}\alpha&\gamma\\ \gamma&\beta\end{smallmatrix}\bigr)\otimes\id$.
With $G=\bigl(\begin{smallmatrix}0&-1\\1&0\end{smallmatrix}\bigr)\otimes\id$ one finds $\tfrac12(N-GNG)=\tfrac12(\alpha+\beta)\id$ and
$\tfrac12(N+GNG)=\tfrac12(\alpha-\beta)\eta-\gamma\,\eta G$. Since $G^{\!\top}G=\id$ and $G^{\!\top}=-G$, $N$ is invariant under $e^{\vartheta G}$ iff
$G^{\!\top}N+NG=0$ iff $[N,G]=0$, which gives (ii). For (iii), the boost law gives
$\delta(\alpha|\bB|^2+\beta|\bE|^2+2\gamma\,\bB\cdot\bE)=-2(\alpha+\beta)\,\bv\cdot(\bE\times\bB)$, the $\gamma$ term being unchanged, so $\beta=-\alpha$.
\end{proof}

The theorem restates \cref{thm:inertiareq}(iii) and \cref{prop:lGhoms} on $V$ rather than on $V^*$, and it is inherited: for every conjugate
root $G'=gGg^{-1}$, $g\in SO(3,3)$, the forms invariant under the compact part of $\mathfrak l_{G'}$ and under $G'$ are the multiples of the
positive form $g^{-\top}g^{-1}$, so the statement could not have failed. It is the uniqueness of the energy density among rotation- and
duality-invariant quadratic forms. What the arena adds is that this form is the Euclidean form of its own basis: $G$ is $\eta$-self-adjoint and
anticommutes with $\Ad_{e_1e_2e_3}|_V=\eta$ (\cref{lem:Gclifford}), so $G^{\!\top}G=\eta G\eta G=\id$. $E_6$ has duality weight zero, while
$e^{\vartheta G}$ turns $(\Re q,\Im q)$ through $2\vartheta$. In this sense $E_6$ is undivided: it distinguishes neither the magnetic triple from the
electric one nor, as \cref{sec:energy-identity} shows, rest from motion.

\begin{proposition}[The pointwise symmetry of the undivided energy]
\label{en:u3}
\begin{enumerate}[label=(\roman*),leftmargin=2.2em]
\item The maps of $V$ that preserve $E_6$ and commute with $G$ form $\uu3$, of dimension $9$: $A=\bigl(\begin{smallmatrix}P&Q\\-Q&P\end{smallmatrix}\bigr)$
acts on $z$ by the anti-Hermitian matrix $P-iQ$.
\item The Lie algebra of $O(3,\C)\times U(1)_G$, the $G$-linear maps preserving $q$ up to a phase, is $\lG\oplus\R G$, of dimension $7$.
\item The intersection of the two is $\so3\oplus\R G$, of dimension $4$, the maximal compact subalgebra of $\lG\oplus\R G$; and $\uu3\cap\lG=\so3$.
\item The stabiliser of the function $E_0$ in $U(3)$ is $U(1)\cdot O(3)=\{e^{i\varphi}O\}$. For instance the unitary map
\[
U=\frac1{\sqrt2}\begin{pmatrix}1&1&0\\ i&-i&0\\ 0&0&\sqrt2\end{pmatrix}
\]
carries the field $(1,0,0)$, at rest with $E_0=E_6=\tfrac12$, to the null field $\tfrac1{\sqrt2}(1,i,0)$ at the same $E_6$.
\end{enumerate}
\end{proposition}

\begin{proof}
(i) Commuting with $G$ forces the block form, which is multiplication by $P-iQ$ on $z$; $A^{\!\top}=-A$ means $P^{\!\top}=-P$ and $Q^{\!\top}=Q$, so
$P-iQ$ is anti-Hermitian, and $3+6=9$. (ii) $q(e^{tM}z)=e^{i\varphi t}q(z)$ gives $M^{\!\top}+M=i\varphi\,\id$, so $M=\tfrac{i\varphi}2+M_0$ with
$M_0\in\so{3,\C}=\lG$ (\cref{thm:GLorentz}), and $i\,\id$ is $G$. (iii) A complex antisymmetric anti-Hermitian matrix is real. (iv) If $U\in U(3)$
preserves $|q|$, the form $z^{\!\top}U^{\!\top}Uz$ vanishes on the cone $q=0$; $q$ is irreducible, so $U^{\!\top}U=e^{i\varphi}\id$, and
$O=e^{-i\varphi/2}U$ is unitary and complex orthogonal, hence real.
\end{proof}

All four parts are inherited Hermitian linear algebra, which makes the word \emph{unitary} literally true of the pointwise energy; part (iv) is
the stabiliser of $|z\cdot z|$ among unitary maps of the Riemann--Silberstein vector \cite{Silberstein1907,BialynickiBirula13}. Read as a
statement about the rest energy, (iv) says exactly in which sense $E_0$ reduces a symmetry: $U(3)$ is reduced to $U(1)\cdot O(3)$, and the
five-dimensional orbit $U(3)/U(1)\cdot O(3)$ parametrises the real structures $z\mapsto\bar z$ up to phase, of which the arena has chosen one
(\cref{en:parity}). But $U(3)$ is a symmetry of the pointwise form $E_6$ only. It is not a symmetry of Maxwell's equations on the frame,
$\partial_tz=-i\curl z$ (\cref{prop:RS}): the map $U$ of (iv) carries the wave $(1,i,0)e^{ik(x_3-t)}$ of \cref{ex:circular} to
$\tfrac{1+i}{\sqrt2}(1,1,0)e^{ik(x_3-t)}$, which does not solve them. Nor is it a symmetry of the split-form interior, since $G=i\,\id$ lies in
$U(3)$ and reverses $\eta$, and $U(3)\cap O(\eta)=O(3)$. Of the closures of \cref{sec:closure} it is a symmetry only for the inertia $N=\id/m$,
because the $U(3)$-invariant symmetric forms are the multiples of $\id$.

\subsection{The rest energy is the complete invariant of Lorentz transformations and duality}
\label{sec:energy-orbits}

\begin{theorem}[Orbits]
\label{en:orbits}
\begin{enumerate}[label=(\roman*),leftmargin=2.2em]
\item Each level set $\{|q|=r^2\}$, $r>0$, is one orbit of $O(3,\C)\times U(1)_G$, of real dimension $5$, with two-dimensional stabilisers; every
non-null $z$ is carried to $(r,0,0)$ with $r=\sqrt{2E_0}$.
\item The nonzero null fields form one orbit of $SO(3,\C)$ alone, of dimension $4$, and at a null point $Gz$ is tangent to the Lorentz orbit.
\item Hence $E_0$ is a complete invariant of Lorentz transformations and duality on non-null fields, and $E_6$ is the complete invariant of
$U(3)$, whose orbits are the spheres.
\end{enumerate}
\end{theorem}

\begin{proof}
(i) Write $q(z)=r^2e^{i\varphi}$ and $w=e^{-i\varphi/2}z/r$, a duality rotation, so that $q(w)=1$. Since $q$ is nondegenerate on
$w^\perp=\{f:f\cdot w=0\}$, some $f\in w^\perp$ has $f\cdot f=1$; with $f_3=w\times f$, the identity
$(a\times b)\cdot(a\times b)=(a\cdot a)(b\cdot b)-(a\cdot b)^2$ gives $f_3\cdot f_3=1$, so $A=(w\,f\,f_3)^{\!\top}\in SO(3,\C)$ and
$Ae^{-i\varphi/2}z=re_1$. An element $M_0+\tfrac{i\varphi}2$ fixes $e_1$ iff $\varphi=0$ and $M_0$ is a complex rotation about $e_1$; the stabiliser is
$\so{2,\C}$ and the orbit has dimension $7-2=5$. (ii) A rotation carries a null $z$ to $\lambda(1,i,0)$ with $\lambda>0$. $K_3$ acts by
$\delta z=i\,e_3\times z$, and $i\,e_3\times(1,i,0)=(1,i,0)$, so the boost $e^{\chi K_3}$ with $\chi=-\log\lambda$ finishes. The Lorentz tangent space
at a null $z$ is its $q$-orthogonal complement, which contains $iz$; at $(1,i,0)$, $iz=-e_3\times z=-L_3z$ is the action of a rotation about $e_3$. (iii)
follows from (i) and (ii).
\end{proof}

This is inherited, the classification of a field strength at a point \cite{Synge56,LandauLifshitz2}. Without duality the complete invariant
on nonzero fields is $q$ itself (the transitivity used in the proof of \cref{thm:achiral}), and duality removes its phase, which is twice the
complexion of \cref{prop:planes}. At the null point $(1,i,0)$ the tangent $Gz=-L_3z$ is generated by the rotation about the propagation axis, the identity of the $G$-charge and the
helicity of \cref{ex:Gcharge}.

\subsection{The differentiation identity}
\label{sec:energy-identity}

The relation $u^2-|\mathbf S|^2=\tfrac14|q|^2$ was used in \cref{sec:orth} for the bound \eqref{eq:EBbound}. Its structure is the following.

\begin{proposition}[Differentiation of the undivided energy]
\label{en:split}
For every $z\in\C^3$
\begin{equation}
\label{en:lagrange}
|z|^4=|z\cdot z|^2+|\bar z\times z|^2,\qquad\text{equivalently}\qquad E_6^2=E_0^2+|\mathbf S|^2 .
\end{equation}
The rank-one Hermitian matrix $H=zz^\dagger$ carries all three terms: $\tr H=2E_6$, $(\Im H)_{jk}=\epsilon_{jkl}S_l$, $\tr(H\bar H)=|q|^2=4E_0^2$ and
$H^2=(\tr H)H$. Hence $E_0\le E_6$, with equality iff $\mathbf S=0$, iff $\bar z\parallel z$, iff $\bB\parallel\bE$ or one of them vanishes.
\end{proposition}

\begin{proof}
$\bar z\times z=-2i\,\bE\times\bB$, and $(a\times b)\cdot(c\times d)=(a\cdot c)(b\cdot d)-(a\cdot d)(b\cdot c)$ at $(\bar z,z,z,\bar z)$ gives
$|\bar z\times z|^2=|z|^4-|q|^2$. Next $H_{jk}=B_jB_k+E_jE_k+i\epsilon_{jkl}(\bE\times\bB)_l$ and $\tr(H\bar H)=|\sum_jz_j^2|^2$. The equality case is the
equality case of the Cauchy--Schwarz inequality $|h(\bar z,z)|\le|z|^2$.
\end{proof}

The identity is inherited, a complex Lagrange identity. It is the exact form of the differentiation asked about: the undivided $E_6$ is the
hypotenuse of a right triangle whose legs are a Lorentz-invariant rest part $E_0$ and a motion part $|\mathbf S|$, and the split needs no
interaction and no second field. In Minkowski terms $E_0$ is the Lorentz length of $(E_6,\mathbf S)$, although $(E_6,\mathbf S_{\rm t})$ is not a
four-vector, as the next proposition makes precise. To avoid a clash with the trivector $e_1e_2e_3$ we write $\Theta^{\mu\nu}$ for the
stress--energy tensor.

\begin{proposition}[Rainich identity]
\label{en:rainich}
For the stress--energy tensor of $(\bE,\Btext)$, $\Theta^{00}=\tfrac12(|\bE|^2+|\Btext|^2)=E_6$, $\Theta^{0i}=(\mathbf S_{\rm t})_i$ and
$\Theta^{ij}=\Theta^{00}\delta_{ij}-E_iE_j-B_{{\rm t}i}B_{{\rm t}j}$, one has $\Theta^\mu{}_\mu=0$ and
\begin{equation}
\label{en:rainicheq}
\Theta^\mu{}_\lambda\Theta^\lambda{}_\nu=E_0^2\,\delta^\mu{}_\nu ,\qquad \Theta_{\mu\nu}\Theta^{\mu\nu}=4E_0^2 .
\end{equation}
$\Theta^{\mu\nu}$ is invariant under duality, and for every unit timelike observer $t$ the density $j_t^\mu=\Theta^\mu{}_\nu t^\nu$ is future
causal with $\eta(j_t,j_t)=E_0^2$.
\end{proposition}

\begin{proof}
Duality invariance is a polynomial identity. Both sides of \eqref{en:rainicheq} are covariant, so by \cref{en:orbits} it suffices to take
$z=(r,0,0)$, where $\Theta^\mu{}_\nu=\tfrac{r^2}2\operatorname{diag}(1,1,-1,-1)$; null fields follow by continuity. That $j_t$ is causal and future
directed is the dominant energy condition of the Maxwell field \cite{HawkingEllis73}.
\end{proof}

This is inherited \cite{Rainich1925,Synge56}. The rest energy of a field point is the Lorentz scalar of the tensor whose $00$ component is the
undivided energy; but each observer $t$ sees a different density $j_t$, all of Lorentz length $E_0$, and no single four-vector carries the
energy and momentum of a field point.

\subsection{The rest energy as a minimum, and two boost factors}
\label{sec:energy-min}

\begin{theorem}[Variational form of the rest energy]
\label{en:min}
\begin{enumerate}[label=(\roman*),leftmargin=2.2em]
\item For non-null $z$, $E_0$ is the minimum of $E_6$ over the Lorentz orbit, and over the orbit of Lorentz transformations and duality,
attained exactly where $\mathbf S=0$; for null $z\neq0$ the infimum $0$ is not attained. Moreover $E_6=E_0/\sqrt{1-v_E^2}$ with $v_E:=|\mathbf S|/E_6<1$.
\item If $\bB\parallel\bE\parallel\mathbf n$ and the boost $e^{\chi K}$ of rapidity $\chi$ has axis orthogonal to $\mathbf n$, then
$E_6=E_0\cosh2\chi$, $|\mathbf S|=E_0\sinh2\chi$ and $v_E=\tanh2\chi=2\beta/(1+\beta^2)$ with $\beta=\tanh\chi$. A boost along $\mathbf n$ changes
neither $E_6$ nor $\mathbf S=0$.
\item The observer moving with velocity $\mathbf v_{\rm rest}=\mathbf S_{\rm t}/(E_6+E_0)$ sees $\mathbf S=0$ and $E_6=E_0$, and
$|\mathbf v_{\rm rest}|=|\mathbf S|/(E_6+E_0)$, which is $\beta$ in the setting of (ii).
\end{enumerate}
\end{theorem}

\begin{proof}
(i) By \eqref{en:lagrange}, $E_6\ge E_0$ with equality iff $\mathbf S=0$, and \cref{en:orbits} supplies points with $\mathbf S=0$ in both orbits; for null
$z$ the boost of \cref{en:orbits}(ii) drives $E_6$ to $0$. Dividing \eqref{en:lagrange} by $E_6^2$ gives the last formula. (ii) Take $\mathbf n=e_1$
and $X=\sigma_1e_1+\sigma_2e_4$. By \cref{ex:emboost} and the same computation for the other basis vectors, $[K_3,e_1]=e_5$, $[K_3,e_5]=e_1$,
$[K_3,e_4]=-e_2$ and $[K_3,e_2]=-e_4$, so $\ad_{K_3}^3=\ad_{K_3}$ and $e^{\chi\ad_{K_3}}=\id+\sinh\chi\,\ad_{K_3}+(\cosh\chi-1)\ad_{K_3}^2$. Hence
$\bB=(\sigma_1\cosh\chi,-\sigma_2\sinh\chi,0)$, $\bE=(\sigma_2\cosh\chi,\sigma_1\sinh\chi,0)$ and $\mathbf S=-E_0\sinh2\chi\,e_3$; likewise for $K_2$,
while $K_1$ fixes $e_1$ and $e_4$. (iii) $|\mathbf S|/(E_6+E_0)=\sinh2\chi/(\cosh2\chi+1)=\tanh\chi$. For the direction, $e^{\chi K_3}$ is the passage to
a frame moving with velocity $-\tanh\chi\,e_3$ (after \cref{thm:emfield}), in which the rest frame moves along $+e_3$, and so does
$\mathbf S_{\rm t}=-\mathbf S$.
\end{proof}

The theorem is inherited \cite{LandauLifshitz2}. Its velocity is the rest-frame velocity of \cref{sec:orth}, where the Poynting vector is the
textbook $\mathbf S_{\rm t}$, as \cref{sec:conventions} stipulates for formulas from the literature; with the frame vector $\mathbf S=\bE\times\bB$
the direction would be reversed. A field point carries two boost factors, $E_6=E_0\cosh2\chi$ against $E=mc^2\cosh\chi$ for a particle, because the
energy density is a component of a second-order tensor and is quadratic in the boost: the rest frame moves with $\beta$, the energy with the
relativistic sum $2\beta/(1+\beta^2)$ of $\beta$ with itself, and an observer riding with the energy velocity sees the same $E_6$ and the reversed
Poynting vector. $E_0$ is the rest-energy density of a field point, not the rest energy of a particle carried by the flux.

\subsection{Duality on the invariants, and the stabilisers of field points}
\label{sec:energy-duality}

The duality rotation $D_\vartheta=e^{\vartheta G}=\cos\vartheta\,\id+\sin\vartheta\,G$ acts as $z\mapsto e^{i\vartheta}z$: it preserves $E_6$ and
$\Theta^{\mu\nu}$, multiplies $q$ by $e^{2i\vartheta}$, preserves the positive form for every $\vartheta$ and $\eta$ only for $\vartheta\in\pi\Z$
(\cref{sec:orth}). For $q\neq0$ exactly the angles $\vartheta\in\pi\Z$ fix $q$, and for $q=0$ every angle does. If $q\neq0$, \cref{prop:planes}
supplies Lorentz coordinates with $z=(f_m+if_e)\hat e_1$, $f_m,f_e$ real; then $q=(f_m^2+f_e^2)e^{2i\varphi}$ with $\varphi=\arg(f_m+if_e)$, and
$D_{-\varphi}$ makes the field purely magnetic, $z=|f_m+if_e|\,\hat e_1$. The angle must be the argument and not the principal value of
$\arctan(f_e/f_m)$, which gives $-|f_m+if_e|\,\hat e_1$ when $f_m<0$. All of this is the inherited duality rotation of the Maxwell invariants
\cite{Rainich1925,Calkin65,DeserTeitelboim76}.

\begin{proposition}[The rest energy is the doubly invariant modulus]
\label{en:doubly}
Of the four functions $E_6$, $\Re q$, $\Im q$ and $E_0$, the first is invariant under $U(1)_G$ but not under $\lG$, the next two are invariant under
$\lG$ but not under $U(1)_G$, and $E_0=\tfrac12|q|$ is invariant under both. On the space of Lorentz orbits of nonzero vectors, labelled by
$q\in\C$, duality rotates the label through $2\vartheta$; the stabiliser of a nonzero label is $\{D_0,D_\pi\}=\{\pm\id\}$.
\end{proposition}

\begin{proof}
The six generators of $\lG$ commute with $G$ and lie in $\so{3,3}$, so $\eta(X,AX)=\eta(X,GAX)=0$; a boost $K_i$ is symmetric for the positive form
and $X^{\!\top}K_iX\neq0$ for generic $X$. The labelling is the transitivity of $\lG$ on the level sets of $q$ with the origin removed (proof of
\cref{thm:achiral}), and $D_\vartheta$ acts on the label by $q\mapsto e^{2i\vartheta}q$.
\end{proof}

The proposition is inherited. With $\mathbf g=\mathbf S_{\rm t}/c^2$ and units restored, \eqref{en:lagrange} reads $u^2-c^2|\mathbf g|^2=E_0^2$, the
field form of $E^2-p^2c^2=m^2c^4$ (\cref{en:SI}). A nonzero rest energy thus breaks $U(1)_G$ to $\Z_2$ on the space of Lorentz orbits, and $\arg q$ is the angle along the broken
direction. That is a
reading of the proposition, and it comes with two qualifications. The breaking is a property of a configuration, not of a vacuum: the energy is
minimised by the zero field and no potential selects $|q|\neq0$ (\cref{sec:orth}). And the broken symmetry is not one of the interior Lagrangian,
which duality rotates into its $\eta G$ partner.

\begin{theorem}[Stabilisers of field points]
\label{en:stab}
Let $z\neq0$.
\begin{enumerate}[label=(\roman*),leftmargin=2.2em]
\item $\Stab_{\lG}(z)=\{(\zeta z)\times\cdot\,:\,\zeta\in\C\}$ is abelian of real dimension two. If $q\neq0$ and $\hat z=z/\sqrt q$, then
$R=\hat z\times\cdot$ satisfies $R^3=-R$ (a rotation) and $(GR)^3=GR$ (a boost) about the common complex axis $\hat z$: the stabiliser is
$\so2\oplus\so{1,1}$. If $q=0$, every element $N$ of the stabiliser has $N^3=0\neq N^2$: null rotations, $\R^2$.
\item In $\lG\oplus\R G$ a non-null $z$ has the same stabiliser, with no $G$ component. A null $z=\mathbf n_1+i\mathbf n_2$ has a three-dimensional
stabiliser $\cong\mathfrak e(2)$, spanned by the null rotations and $L_{\mathbf n}+G$, with $\mathbf n=\mathbf n_1\times\mathbf n_2/|\mathbf n_1|^2=\bB\times\bE/|\bB\times\bE|$
the propagation axis, which acts on the null rotations with eigenvalues $\pm i$. For $z=(1,i,0)$, $(L_3+G)z=0$ and $e^{\vartheta G}z=e^{-\vartheta L_3}z$.
\item In $\lint$, on the interior four-plane $\langle e_1,e_4,e_5,e_6\rangle$, $\Stab(e_1)=\so3$ and $\Stab(e_1+e_6)=\mathfrak e(2)$, both
three-dimensional.
\end{enumerate}
\end{theorem}

\begin{proof}
(i) $w\times z=0$ iff $w\in\C z$. From $(z\times)^2w=z(z\cdot w)-qw$ one has $(z\times)^3=-q\,(z\times)$; with $\hat z\cdot\hat z=1$ this is $R^3=-R$, and
$G\leftrightarrow i$ gives $(GR)^3=GR$; for $q=0$, $(z\times)^3=0\neq(z\times)^2$. (ii) The condition is $w\times z+i\vartheta z=0$, and dotting with $z$ gives
$i\vartheta q=0$. For null $z$, $\mathbf n_1\cdot\mathbf n_2=0$ and $|\mathbf n_1|=|\mathbf n_2|$ give $\mathbf n\times\mathbf n_1=\mathbf n_2$ and
$\mathbf n\times\mathbf n_2=-\mathbf n_1$, so $\mathbf n\times z=-iz$, $(L_{\mathbf n}+G)z=0$ and $[L_{\mathbf n},z\times\cdot]=(\mathbf n\times z)\times\cdot=-i(z\times\cdot)$.
(iii) is the classical computation of the little groups of a timelike and a null vector.
\end{proof}

The theorem is inherited: (iii) is Wigner's classification \cite{Wigner39,WeinbergQFT}, and (i)--(ii) are the classical stabilisers of a
bivector and the identity of helicity and duality on null fields \cite{PenroseRindler84,Calkin65}, which on the frame are \cref{thm:axis} (the
axis of a null field) and \cref{ex:Gcharge} (the $G$-charge is the helicity). It separates two notions that the word \emph{mass} runs together.
The little group of a momentum is three-dimensional with or without mass, and the stabiliser of a field point is two-dimensional with or without
$q$. A mass changes the type of a Lorentz stabiliser, never its dimension, and breaks no Lorentz generator, since $q$ and the square of a momentum
are invariants. What a nonzero $q$ removes is the duality generator: a null field point is invariant under duality up to a rotation about its
axis, so the group $\mathfrak e(2)$ of Wigner reappears once duality is admitted, with the helicity as its compact generator, while for a non-null
point duality is broken to $\{\pm\id\}$.

\subsection{\texorpdfstring{Where $c$ enters}{Where c enters}}
\label{sec:energy-c}

Adding the dimensionless $G$ to a derivative along a clock coordinate needs a length $\ell$ and gives the rate $\kappa=c/\ell$ of
\eqref{eq:kappa}. On the frame, $c$ enters as follows.

\begin{proposition}[$c^2$ is an isotropic weight]
\label{en:onec}
Let the temporal triple carry $\bB$ and the spatial triple $\bE$, let $C:V^+\to V^-$ be the conversion and
$G_C=\bigl(\begin{smallmatrix}0&-C^{-1}\\ C&0\end{smallmatrix}\bigr)$, so that $G_C^2=-\id$.
\begin{enumerate}[label=(\roman*),leftmargin=2.2em]
\item $G_C$ commutes with the diagonal rotations iff $C=c\,\id$.
\item For $C=c\,\id$, writing $G_c$ for $G_C$, the forms invariant under the diagonal rotations and under $e^{\vartheta G_c}$ are the multiples of $c^2|\bB|^2+|\bE|^2$.
\item $(\kappa G)^2=-c^2/\ell^2\,\id$, and $\ell$ is not fixed.
\end{enumerate}
\end{proposition}

\begin{proof}
(i) The lower-left block of $[G_C,R\oplus R]$ is $CR-RC$, and Schur's lemma gives $C=c\,\id$; this is \cref{thm:onec}. (ii) For
$N=\bigl(\begin{smallmatrix}\alpha&\gamma\\ \gamma&\beta\end{smallmatrix}\bigr)\otimes\id$,
$G_c^{\!\top}N+NG_c=\bigl(\begin{smallmatrix}2c\gamma&c\beta-\alpha/c\\ c\beta-\alpha/c&-2\gamma/c\end{smallmatrix}\bigr)\otimes\id$, which vanishes iff
$\gamma=0$ and $\alpha=c^2\beta$. (iii) $G^2=-\id$.
\end{proof}

Part (i) is \cref{thm:onec}, fixed there, and it fails for $C=\operatorname{diag}(c_1,c_2,c_3)$ with distinct entries; part (ii) is inherited, and
part (iii) is constructed, a definition, since $G^2=-\id$ is adopted in \cref{thm:neutral} and $\ell$ is chosen. Even the weight is a unit convention: $\Phi=\operatorname{diag}(c,c,c,1,1,1)$
conjugates $G_c$ to $G$ and turns $c^2|\bB|^2+|\bE|^2$ into $|\Phi X|^2$, so with $\bB$ measured as $c\bB$, as in Heaviside--Lorentz units, $c$
disappears from $E_6$ and appears only in the rest-mass density $\rho_0$ below.

\begin{proposition}[SI form]
\label{en:SI}
With $\varepsilon_0\mu_0c^2=1$, $z_c:=c\bB+i\bE$, $u=\tfrac{\varepsilon_0}2|z_c|^2$, $\mathbf S_{\rm t}=\bE\times\Btext/\mu_0$, $\mathbf g=\mathbf S_{\rm t}/c^2$,
$q_c:=z_c\cdot z_c=c^2|\Btext|^2-|\bE|^2-2ic\,\Btext\cdot\bE$ and $u_0:=\tfrac{\varepsilon_0}2|q_c|$, \cref{en:split} applied to $z_c$ gives
\begin{equation}
\label{en:SIeq}
u^2=c^2|\mathbf g|^2+u_0^2,\qquad \rho_0:=u_0/c^2,\qquad [u]=[c\,\mathbf g]=[u_0]=\mathrm{J\,m^{-3}},\quad [\rho_0]=\mathrm{kg\,m^{-3}} .
\end{equation}
\end{proposition}

This is inherited, and calling $\rho_0$ a rest-mass density is a reading, since $(u,c\,\mathbf g)$ is not a four-vector (\cref{en:rainich}).
It is the exact content of the clause, taken up in \cref{sec:energy-readings}, that $c^2$ is the symmetric part of $E=mc^2$, and it is less than the statement suggests. The constant $c$
is the unique conversion between the triples that commutes with their common rotations, which are the residual symmetry of both $E_6$ and $E_0$;
in SI units it enters $E_6$ as the isotropic weight $c^2$ of $|\bB|^2$, which a choice of units removes. It is a conversion, not a symmetry, and
nothing breaks it: $u_0$ uses the same $c$, and $\rho_0=u_0/c^2$ only converts units.

\subsection{The mass of a whole}
\label{sec:energy-whole}

For a free field with compactly supported Cauchy data, $(U,\mathbf P)=\int(u,\mathbf S_{\rm t})\,d^3x$ is a four-vector and its Lorentz length $M$
is positive whenever $q\neq0$ on a set of positive measure, while $M$ does not detect orthogonality (\cref{sec:orth}). The following bounds and
examples locate $M$ relative to the rest energies of the field points.

\begin{proposition}[Mass of a whole]
\label{en:whole}
\begin{enumerate}[label=(\roman*),leftmargin=2.2em]
\item $M^2=U^2-|\mathbf P|^2\ge U\int(u-|\mathbf S|)\,d^3x\ge0$, and in every frame $M\ge\int E_0\,d^3x$, with equality iff the densities
$(u,\mathbf S_{\rm t})(x)$ are all proportional.
\item A plane wave has $M=0$. In SI units the waves $E_x=a\cos(kz\mp\omega t)$, $cB_{{\rm t}y}=\pm a\cos(kz\mp\omega t)$ superpose to
$E_x=2a\cos kz\cos\omega t$, $cB_{{\rm t}y}=2a\sin kz\sin\omega t$. Per cell (unit area, length $2\pi/k$) one has $U=2\pi\varepsilon_0a^2/k$ and
$\mathbf P=0$, so $Mc^2=U$, although each constituent has $q_c=0$ and $U_\pm=|\mathbf P_\pm|c$. The total has $q_c=-2a^2(\cos2kz+\cos2\omega t)$,
$\int_{\rm cell}\Theta^{zz}=U$ and $\int_{\rm cell}\Theta^{xx}=-U\cos2\omega t$, and the integrated rest energy $\int_{\rm cell}u_0$ oscillates between
$\tfrac2\pi U$ and $U$.
\item With $c=1$ and $f(s)=(1-s^2)^2$ for $|s|\le1$, $f(s)=0$ otherwise, the field $E_x=f(z-t)+f(z+t)$, $B_{{\rm t}y}=f(z-t)-f(z+t)$ has, per unit
transverse area, $U=\tfrac{512}{315}$ and $\mathbf P=0$ at all times, so $M=U$, with $(U,P_z)$ a two-vector under boosts along $z$ only. For $|t|\ge1$
the two pulses are disjoint, $q\equiv0$ and $\int E_0\,d^3x=0$.
\end{enumerate}
\end{proposition}

\begin{proof}
(i) Conservation and finite propagation speed make $\int_\Sigma\Theta^{\mu\nu}d\Sigma_\nu$ independent of the spacelike hyperplane $\Sigma$
\cite{Laue1911}. Since $u-|\mathbf S|=E_0^2/(u+|\mathbf S|)\ge0$ and $M^2=(U-|\mathbf P|)(U+|\mathbf P|)$, the first bound follows. Each $(u,\mathbf S_{\rm t})(x)$ is
future causal of Lorentz length $E_0(x)$ (\cref{en:rainich}), and future causal vectors obey the reversed triangle inequality
$|P_1+P_2|\ge|P_1|+|P_2|$, with equality iff they are proportional. (ii) The cell integrals are elementary. For the last statement, with
$s=\sin^2\omega t$ one has $q_c=4a^2(s-\cos^2kz)$, so $u_0=2\varepsilon_0a^2|s-\cos^2kz|$; its cell mean is least at the median $s=\tfrac12$ of the
arcsine distribution of $\cos^2kz$, where it is $\tfrac2\pi U$, and equals $U$ at $s=0$ and $s=1$. (iii) $u=f(z-t)^2+f(z+t)^2$ and $S_{{\rm t}z}=f(z-t)^2-f(z+t)^2$, the
cross terms cancelling, and $q=-4f(z-t)f(z+t)$; finally $2\int_{-1}^1(1-s^2)^4ds=\tfrac{512}{315}$.
\end{proof}

The first bound and item (ii) are inherited. The cell in (ii) is not isolated: $\int\Theta^{zz}=U$ is the pressure on whatever confines it, so its
$(U,\mathbf P)$ is a four-vector only with the walls included, as for radiation in a box \cite{Laue1911}. The inequality $M\ge\int E_0$ and
item (iii) obstruct the reading that the mass of a whole emerges from the interference that makes $q\neq0$: interference suffices for $M>0$, by (i) and the example (ii),
but it is not necessary, by (iii). A whole gets its mass from the non-collinearity of the energy--momenta of its parts, as a pair of photons does.

\subsection{Splitting the temporal energy lowers the mass}
\label{sec:energy-temporal}

There is a second reading of the phrase that the energy is six-dimensional: the three temporal directions of $V^+$ carry three energies, and a
mass would arise by differentiating them. On the arena this reading fails.

\begin{proposition}[Splitting the temporal energy lowers the mass]
\label{en:temporal}
Let $P=\sum_i(\mathsf E_i/c)e_i+\sum_ap_ae_{a+3}$ with $\eta(P,P)=\mu^2c^2$, let $\widehat W$ be the symmetric clock of \cref{prop:clock},
$\mathsf E_W:=(\mathsf E_1+\mathsf E_2+\mathsf E_3)/\sqrt3$, and let $\mathsf E_\perp$ be the part of $(\mathsf E_1,\mathsf E_2,\mathsf E_3)$ orthogonal
to $(1,1,1)$, the energies of the relative times. The mass read by the clock, $m_4^2c^4:=\mathsf E_W^2-|\bp|^2c^2$, is
\[
m_4^2c^4=\mu^2c^4-|\mathsf E_\perp|^2\le\mu^2c^4 .
\]
For the clock $e_1$ instead, $m_4^2c^4=\mu^2c^4-\mathsf E_2^2-\mathsf E_3^2$; for example $\mathsf E=(3,4,0)$, $\bp=(5,0,0)$ and $c=1$ give $\mu=0$ and
$m_4^2=-16$.
\end{proposition}

\begin{proof}
$|(\mathsf E_1,\mathsf E_2,\mathsf E_3)|^2=\mathsf E_W^2+|\mathsf E_\perp|^2$ and $\mu^2c^4=c^2\eta(P,P)=|(\mathsf E_1,\mathsf E_2,\mathsf E_3)|^2-|\bp|^2c^2$.
\end{proof}

This is an obstruction, and it restates \cref{prop:toy} and \cref{ex:tachyons} in a form covariant under the choice of clock. In the momentum
reading the differentiation of the temporal energy subtracts from the mass and can make it tachyonic, and under the surplus constraints of
\cref{thm:wlgauge}, $\mathsf E_\perp=0$, nothing is left to differentiate. In the field-strength reading the motion part adds instead,
$E_6\ge E_0$. Only the latter carries the structure of an undivided energy with an internal differentiation.

\subsection{\texorpdfstring{Three candidates for $m$ in the arena}{Three candidates for m in the arena}}
\label{sec:energy-candidates}

The arena has three structures that might be called a mass: the temporal volume element $T=e_1e_2e_3$, which is the Dirac mass form
(\cref{cor:fullmass}); the parameter $m$ and the conserved energy of the interior dynamics of \cref{sec:neutral}; and the rest energy of the
momentum field in the closure of \cref{sec:closure}. We take them in turn.

\begin{proposition}[The mass form acts on the frame as parity]
\label{en:parity}
\begin{enumerate}[label=(\roman*),leftmargin=2.2em]
\item $\Ad_T|_V=\eta$ fixes $\bB$ and reverses $\bE$: on the frame it is parity. It commutes with the $L_i$ and anticommutes with $G$ and the
$K_i$, so it normalises $\lG$ without being $\lG$-invariant.
\item $\eta X$ corresponds to $\bar z$ and $q(z)=h(\bar z,z)$, so $E_6=\tfrac12h(z,z)$ and
$E_0=\tfrac12|h(\bar z,z)|=\tfrac12\max_\vartheta\eta(e^{\vartheta G}X,e^{\vartheta G}X)$.
\item As a Grassmann mass kernel on spinors, $T$ is invariant under all of $\so{3,3}$, $\lG$ included, in the sense of \cref{def:subkernel}, and
breaks nothing; only its sign changes, under conjugation by a spatial generator.
\end{enumerate}
\end{proposition}

\begin{proof}
(i) $\Ad_T|_V=\eta$ is \cref{lem:Gclifford}; $\eta G\eta=-G$ blockwise, the $L_i$ preserve each triple and the $K_i$ exchange them. (ii)
$z(\eta X)=\bB-i\bE$, $h(\bar z,z)=\sum_jz_j^2$, and $\eta(e^{\vartheta G}X,e^{\vartheta G}X)=\Re(e^{2i\vartheta}q)$. (iii) is \cref{cor:fullmass}.
\end{proof}

The identities are inherited or restate the paper; their consequence is an obstruction by type. The identification of $m$ with $T$ fails on frame
vectors, where the adjoint action of $T$ is parity, while as a kernel on spinors $T$ is invariant under $\lG$. The Dirac mass lives on the interior
reading of the six directions, which \cref{prop:emexclusive} excludes on the same six directions, and even there it breaks nothing. What survives
is a reading: the rest energy is the overlap $\tfrac12|h(\bar z,z)|$ of a field with its mirror image, and $\Ad_T$, the arena's splitting into
two triples, is the real structure whose stabiliser in $U(3)$ is $U(1)\cdot O(3)$ (\cref{en:u3}(iv)).

The second candidate is the interior. By \cref{prop:orbitsEM}, every solution of $\ddot X=-\kappa^2X$ is $z(\tau)=c_-e^{-i\kappa\tau}+c_+e^{i\kappa\tau}$,
$c_\pm\in\C^3$, the momentum field $\pi=m\alpha_+X$ corresponds to $2im\kappa c_+e^{i\kappa\tau}$, the Noether energy is $\mathcal E=2m\kappa^2\Re(c_-\cdot c_+)$, and the
complex quantity $\mathcal E_\C:=\tfrac m2\bigl(q(\dot z)+\kappa^2q(z)\bigr)=2m\kappa^2\,c_-\cdot c_+$ is conserved. The pointwise energies of the frame
field and of the momentum field behave differently along these solutions.

\begin{proposition}[Frame field and momentum field in the interior]
\label{en:interior}
Along every solution of the split-form interior with the binding root:
\begin{enumerate}[label=(\roman*),leftmargin=2.2em]
\item The momentum field is $z(\pi)=2im\kappa\,c_+e^{i\kappa\tau}$, that is $\pi(\tau)=e^{\kappa\tau G}\pi_0$: it moves by pure duality, and
$E_6(\pi)=\tfrac12|p|^2=2m^2\kappa^2|c_+|^2$, $E_0(\pi)=2m^2\kappa^2|c_+\cdot c_+|$ and $\mathbf S(\pi)$ are constant.
\item The frame field has
\[
|z|^2=|c_-|^2+|c_+|^2+2\Re\bigl(h(c_-,c_+)e^{2i\kappa\tau}\bigr),\qquad q(z)=e^{-2i\kappa\tau}c_-\cdot c_-+2c_-\cdot c_++e^{2i\kappa\tau}c_+\cdot c_+ ,
\]
so its $E_6$ is constant iff $h(c_-,c_+)=0$, and its $E_0$ oscillates in general: $c_-=e_1$, $c_+=e_2$ give $\mathcal E=0$, $E_6=1$ and
$E_0=|\cos2\kappa\tau|$.
\item Both pure senses, $z=e^{-i\kappa\tau}c_-$ and $z=e^{i\kappa\tau}c_+$, are duality rotations of the frame field, with $\mathcal E_\C=0$ and with
constant $E_6$ and $E_0$ of the frame field. On the first, the orbits without momentum, $p=0$ while $E_6=\tfrac12|c_-|^2>0$; on the second
$|p|^2=8m^2\kappa^2E_6$.
\item The parameter $m$ is an overall factor of the Lagrangian and does not enter $\ddot X=-\kappa^2X$.
\end{enumerate}
\end{proposition}

\begin{proof}
$\dot z+i\kappa z=2i\kappa c_+e^{i\kappa\tau}$ is $z(\alpha_+X)$, since $G$ acts as $i$, and $\pi=m\alpha_+X$; with $|\pi|=|p|$ ($\eta$ is orthogonal)
this gives (i), and $\overline{z(\pi)}\times z(\pi)$, $|z(\pi)|$ and $|z(\pi)\cdot z(\pi)|$ are invariant under multiplication of $z(\pi)$ by a phase. (ii) is the
expansion of $\bar z\cdot z$ and $z\cdot z$; in the example $|z|^2=2$ and $q=2\cos2\kappa\tau$. (iii) is (i) and (ii) at $c_+=0$ or $c_-=0$, with
$|p|^2=|z(\pi)|^2=4m^2\kappa^2|c_+|^2$. (iv) is read off \eqref{eq:splitL} and \cref{thm:neutral}.
\end{proof}

Item (i) and the first half of (iii) make explicit what \cref{prop:orbitsEM} states about the momentum and about the orbits without momentum;
item (ii) and the second half of (iii) are fixed by the algebra and complete the picture, in which parity exchanges the two senses
(\cref{thm:EBsplit}). The
quantities the interior conserves belong to the momentum field: $E_6(\pi)$, $E_0(\pi)$ and $\mathbf S(\pi)$ are constant already in the clamped
interior, because $\pi$ moves by a duality rotation. The same functions of the frame field are not conserved, and $E_6$ of the frame field is not a
function of $|p|^2$: on the orbits without momentum $|p|^2=0$ while $E_6>0$, on the duality rotations of the other sense $|p|^2=8m^2\kappa^2E_6$.
Duality acts along both pure senses, so duality orbits are not the orbits without momentum; that sense is singled out by the momentum, not by
duality. And the interior's conserved energy is not undivided: it is the real part of a constant of duality charge two, carried by neither sense
alone.

As an identification of a rest energy, the interior is obstructed. In it $m$ is an overall factor and an input independent of $\kappa$, the pure
duality motion of either sense carries zero energy, and, with $X$ a length and $\kappa=c/\ell$, the energy
$\mathcal E=mc^2\cdot2\Re(c_-\cdot c_+)/\ell^2$ has the form $mc^2$ times a dimensionless amplitude; but the amplitude is indefinite, vanishes on the
pure duality orbits and is selected by nothing, since $\ell$ is free. That $\kappa G$ lies in the $\eta$-self-adjoint summand of
\cref{thm:EAdichotomy} is fixed; nothing in the interior turns $\kappa^2$ into a rest energy, and reading the rate as a mass, $mc^2=\hbar\kappa$, would
be an added identification, not a result.

The third candidate is the closure.

\begin{remark}[The momentum field in the closure]
\label{en:closure}
The closure of \cref{sec:closure}, whose bearing on the rest energy is the subject of \cref{sec:restclosure}, settles the third candidate by
results fixed there. For every Hamiltonian, $E_6(\pi)=\tfrac12|p|^2$, $E_0(\pi)=\tfrac12|w(p)|$, with $w$ as in \eqref{eq:wdef}, and
$\mathbf S(\pi)$ are Casimirs (\cref{prop:LP,cl:frozen}), while the same functions of the frame field oscillate in the clamped interior
(\cref{en:interior}) and, with a free rotor in the rotating regime, Lorentz-class inertia and $E_0(\pi_0)>0$, grow without bound
(\cref{thm:freedrift,cl:nogen}). For Lorentz-class $N=(\nu_1\eta^{-1}+\nu_2G\eta^{-1})/m$ the potential of the rotor angle, which is the
duality angle of the momentum field, has the amplitude $|\nu_1+i\nu_2|\,E_0(\pi_0)/m$ (\cref{thm:pendulum}); for $(\nu_1,\nu_2)\neq(0,0)$ the
rate is constant and nonzero exactly when $E_0(\pi_0)=0$ and the initial rate is nonzero (\cref{cor:constrate}), and if $E_0(\pi_0)>0$ it is
constant only at the equilibria, where it vanishes (\cref{cl:flat}). For the duality-class $N=\id/m$ the frame energy
$\tfrac12p^{\!\top}Np$ is $E_6(\pi)/m$, the undivided energy of the momentum field, and the rate is conserved for all data
(\cref{thm:inertiareq}(iii)).
\end{remark}

Read through a dictionary that is not a theorem, in which the clamped rate $\kappa=c/\ell$, with $\ell$ free, stands for the conversion and
$E_0(\pi)$, which has the dimension of a squared momentum, stands for a rest energy, these results say the following. The momentum field's rate keeps its symmetry, constancy and invariance under the duality translation of the rotor, when the momentum
field is null; its rest part $E_0(\pi)$ is the amplitude that breaks it; and $E_6(\pi)$ is a Casimir and breaks nothing. But the breaking is explicit,
by an inertia that \cref{thm:inertiareq} forbids from also conserving the rate; what is broken is the constancy of the rotor rate, a property of
duality, not Lorentz invariance and nothing carried by $c$; and with the undivided energy as Hamiltonian, $N=\id/m$, the rate is constant at the
price of the drift of the frame in \cref{thm:drift}.

\subsection{What transport by a compact gauge field does to the two energies}
\label{sec:energy-transport}

An admissible Yang--Mills generator is compact: by \cref{thm:sig} and \cref{cor:noboost} gauge generators lie in $\so3_T\oplus\so3_S\subset\so{3,3}$,
the rotations of the two triples. Write such a generator as $b=(A+D,A-D)$, with $A+D$ acting on the temporal and $A-D$ on the spatial triple. Its
diagonal part $(A,A)$ lies in $\lG$ and its anti-diagonal part $(D,-D)$ in $\mathfrak p_G$. Boost-valued connections are legitimate metric
connections (\cref{prop:tear}) but not admissible gauge potentials.

\begin{proposition}[What transport preserves]
\label{en:transport}
\begin{enumerate}[label=(\roman*),leftmargin=2.2em]
\item An element $Y\in\so{3,3}$ preserves $\bE\cdot\bB$ for all $X$, equivalently $q$, equivalently $E_0$, iff $[Y,G]=0$, that is iff $Y\in\lG$.
\item A compact $b=(A+D,A-D)$ preserves $E_6$ and $\Re q=|\bB|^2-|\bE|^2$, and it preserves $E_0$ for all $X$ iff $D=0$. The anti-diagonal
rotation through $\tfrac\pi4$ about $e_3$ carries the null field $(1,i,0)$ to the field $\tfrac{1+i}{\sqrt2}(1,1,0)$ at rest, $E_0$ going from $0$ to $1$
at $E_6=1$.
\item At three constituents the part of the Pati--Salam algebra of $P_*$ commuting with $G_3$, the algebra $\so4_\Delta\cong\su2\oplus\su2$ of
\cref{thm:PSbinding}, consists of paired rotations, so it preserves the $\so{9,\C}$ invariants $\tfrac12|Z|^2=\tfrac12\sum_k|z^{(k)}|^2$ and
$q_3(Z)=\sum_kz^{(k)}\cdot z^{(k)}$, hence $\tfrac12|q_3|$.
\item It preserves neither the rest energies of the constituents nor that of their summed density: the rotation through $\tfrac\pi4$ generated by
$\tfrac12\mathrm{rot}(7,13)\in\so4_\Delta$ maps $(z^{(1)},z^{(2)},z^{(3)})=(0,(1,i,0),(1,0,0))$ to $(0,(0,i,0),(\sqrt2,0,0))$, keeping
$\tfrac12|Z|^2=\tfrac32$ and $q_3=1$, while $\sum_kE_0^{(k)}$ goes from $\tfrac12$ to $\tfrac32$ and
$\bigl(\sum_kE_6^{(k)};\sum_k\mathbf S_{\rm t}^{(k)}\bigr)$ goes from $(\tfrac32;0,0,1)$ to $(\tfrac32;0,0,0)$, its Lorentz length from $\tfrac{\sqrt5}2$
to $\tfrac32$.
\end{enumerate}
\end{proposition}

\begin{proof}
(i) For $Y\in\so{3,3}$, $Y^{\!\top}\eta G+\eta GY=\eta[G,Y]$, so $Y$ preserves the form $\eta G$, that is $\bE\cdot\bB$, iff $[Y,G]=0$. Every $Y\in\so{3,3}$
preserves $\Re q$; if it also preserves $|q|$, then $\Im q\cdot\delta(\Im q)=0$, so $\delta(\Im q)=0$ because polynomials form an integral domain. (ii)
$b$ rotates each triple orthogonally, so it preserves $|\bB|^2$ and $|\bE|^2$; $(A,A)$ commutes with $G$ and $G(D,-D)G^{-1}=(-D,D)$, so by (i) $E_0$ is
preserved for all $X$ iff $D=0$. In the example $\bB=(1,0,0)\mapsto\tfrac1{\sqrt2}(1,1,0)$ and $\bE=(0,1,0)\mapsto\tfrac1{\sqrt2}(1,1,0)$. (iii) By
\cref{thm:so9C} the paired rotations are the elements $\bigl(\begin{smallmatrix}R&0\\0&R\end{smallmatrix}\bigr)$ with $R$ real antisymmetric, which act on
$Z\in\C^9$ by a real orthogonal matrix and preserve $|Z|^2$ and $q_3$; by \cref{thm:PSbinding} $\so4_\Delta$ is spanned by such elements. (iv)
$\mathrm{rot}(7,13)$ rotates the first components of $z^{(2)}$ and $z^{(3)}$ into each other by a real rotation, and the stated values follow by
direct evaluation, with $\mathbf S_{\rm t}=(0,0,1)$ for $(1,i,0)$.
\end{proof}

Parts (i) and (ii) are inherited linear algebra, and they cut both ways. Pointwise, the anti-diagonal part of an admissible $b$ does turn a null
field into one at rest at fixed undivided energy, while the diagonal part, a rotation in $\lG$, changes no invariant. But the anti-diagonal part is
the part that fails to commute with $G$: along a worldline it cannot be removed by a rotating frame, unlike the diagonal part (\cref{bf:remove}; Larmor's theorem
\cite{Larmor1897}), and in the interior it does not commute with the shift $\kappa G$. There, alone, it produces Jordan blocks and secular growth
(\cref{bf:pG}), and a mixed compact connection keeps every orbit bounded exactly when it rotates the two triples at equal rates about axes that
are not opposite, outside one exceptional ratio of the rates to $\kappa$ at which some solutions drift (\cref{bf:screen}); these are statements
about the interior transport, not about the pointwise invariants, and they are not used here. Nothing in the proposition selects an amount of rest energy. Parts (iii) and (iv) are an
obstruction for the $\so{9,\C}$ invariants only: no transport valued in the paired rotations changes $\tfrac12|Z|^2$ or $q_3$. They are no obstruction
for the rest energy of the whole in the diagonal Lorentz reading, which the example changes; $\mathrm{rot}(7,13)$ does not commute with the diagonal
rotations, and which of the two quantities is physical depends on the Lorentz action at three constituents, which the algebra leaves open
(\cref{tab:choices}; by \cref{thm:nocommon} the closure symmetry meets the interior Lorentz algebra of the first constituent in $0$). Moreover
$\so4_\Delta$ is the part of the Pati--Salam algebra commuting with $G_3$, not the only part whose transport keeps orbits bounded (\cref{bf:so4spec}(v)). Here $b$ enters
only as a constant transport element or a holonomy, with no field strength and no energy of its own, so where a gauge field belongs in the
differentiation of the energy is, at this level, a matter of reading.

\subsection{\texorpdfstring{What the arena says about $E=mc^2$}{What the arena says about E=mc2}}
\label{sec:energy-readings}

The results of this section can be held against three clauses that are often attached to $E=mc^2$: that the energy is one undivided,
six-dimensional quantity; that mass arises from an internal differentiation of it; and that $c^2$ is the symmetric part of the relation.

\paragraph{The energy is one undivided, six-dimensional quantity.} This holds in the field-strength reading. Among quadratic forms on the six
frame components exactly one is invariant under the diagonal rotations and under duality, and it is positive: $E_6=\tfrac12(|\bB|^2+|\bE|^2)$
(\cref{en:unique}, inherited, restating \cref{thm:inertiareq}(iii)). Its symmetry group is $U(3)$, so the energy is unitary in a literal sense, but
only as a pointwise quadratic form: $U(3)$ is not a symmetry of Maxwell's equations on the frame nor of the split-form interior (\cref{en:u3}). In
the closure the undivided energy of the momentum field, $E_6(\pi)=\tfrac12|p|^2$, is a Casimir, conserved by every Hamiltonian together with
$E_0(\pi)$ and $\mathbf S(\pi)$ (\cref{en:closure}, restating \cref{prop:LP}); the energies of the frame field are not conserved
(\cref{en:interior,en:closure}). The other reading of the clause, three temporal energies forming a vector in $V^+$, is obstructed: splitting that
vector lowers the mass (\cref{en:temporal}).

\paragraph{Mass arises from an internal differentiation of that quantity.} This holds as an identity, not as a process. With no interaction and no
second field the undivided energy of a field point splits as $E_6^2=E_0^2+|\mathbf S|^2$, $E_0=\tfrac12|z\cdot z|$ (\cref{en:split}), in SI units
$u^2=c^2|\mathbf g|^2+u_0^2$, the field-point form of $E^2=(pc)^2+(mc^2)^2$ (\cref{en:SI}). The rest part is the complete invariant of Lorentz
transformations and duality (\cref{en:orbits}), the Lorentz scalar of the stress tensor whose $00$ component is $E_6$ (\cref{en:rainich}), and the
minimum of $E_6$ over boosts, attained exactly in the rest frame (\cref{en:min}). All of this is inherited classical electrodynamics in the arena's
variables. The arena adds which form plays which role: $E_6$ belongs to the duality class of \cref{thm:inertiareq} and $E_0$, the modulus of the two
Lorentz-class invariants, is built from the Lorentz class, and reading $E=mc^2$ as the locus where the two coincide, the rest frame, is a reading. The
differentiation does not happen in time: in the closure the split of the momentum field is fixed by the initial data, all three terms being
Casimirs (\cref{en:closure}). Nor does a whole get its mass from interference: two disjoint counter-propagating null pulses have $q\equiv0$ everywhere
and $M=U>0$ (\cref{en:whole}).

\paragraph{$c^2$ is the symmetric part.} This holds in one precise form. The constant $c$ is the unique isotropic conversion between the temporal and
the spatial triple, fixed by Schur's lemma (\cref{thm:onec}); it is a conversion, not a symmetry, and its weight in $E_6$ is a unit convention
(\cref{en:onec}). The form in which the clause is exact is the decomposition of the shift into the two summands of \cref{thm:EAdichotomy},
\[
\End(V)=\so{3,3}\oplus\{Y:Y^\dagger=Y\},\qquad \kappa G\in\{Y:Y^\dagger=Y\},\qquad (\kappa G)^2=-\frac{c^2}{\ell^2}\,\id .
\]
The $\eta$-self-adjoint summand is the one in which factorisation holds, the one that binds in neutral signature (\cref{thm:neutral}), and the one
that carries $c$, through $\kappa=c/\ell$; the antisymmetric summand holds the connections. The decomposition and the fact that only the first summand
binds are fixed; the value $(\kappa G)^2=-c^2/\ell^2$ is constructed, a definition with $\ell$ free; and reading the first summand as ``$c^2$'' and the second, which holds the
connections, as the place of a gauge field is a reading; that a gauge generator differentiates the energy is not supported, since it selects no
amount of rest energy (\cref{en:transport}). Its limits are fixed results: duality is a symmetry of the positive form,
while the interior Lagrangian is built on $\eta$, and the two forms lie in the duality and the Lorentz class of \cref{thm:inertiareq}, which meet in $0$;
in the interior $c^2$ enters as $\kappa^2=c^2/\ell^2$ and the energy is $mc^2$ times an indefinite, unselected amplitude (\cref{en:interior}); and in the
closure what $E_0(\pi)$ breaks is the constancy of the rotor rate, not anything carried by $c$ (\cref{en:closure}).

Collected, the clauses become the following statement, which is a reading of fixed and inherited results. The undivided energy is
$E_6=\tfrac12|z|^2$, the one rotation- and duality-invariant quadratic energy of a six-component frame vector, with pointwise symmetry $U(3)$. It
differentiates internally, with no interaction, as $E_6^2=E_0^2+|\mathbf S|^2$, into a Lorentz-invariant rest energy $E_0=\tfrac12|z\cdot z|$ and a
motion part $|\mathbf S|$. What differentiates is the arena's splitting into a temporal and a spatial triple, the real structure $\Ad_{e_1e_2e_3}$,
which reduces $U(3)$ to $U(1)\cdot O(3)$, a symmetry of the pointwise form and not of any dynamics; nothing Lorentzian is broken, and the duality
$U(1)_G$ is broken to $\Z_2$ only on the space of Lorentz orbits and only by a configuration. The constant $c$ is a conversion, not a symmetry, and
$E_0/c^2$ is a rest-mass density of a field point. The statement has three limits: a field point is not a particle (two boost factors, and the mass of
a whole is not interference); in the closure only the momentum field's energies are conserved, while the frame field's oscillate or grow; and a compact
gauge generator selects no amount of rest energy, its anti-diagonal part moving a field between null and rest at fixed $E_6$ and, at three
constituents, $\so4_\Delta$ changing the rest energies of the constituents at fixed $\tfrac12|Z|^2$ and $q_3$. Whether some dynamics selects the rest
frame of a field or the value of $E_0$ is open; nothing in this section does.
